
\noindent
\textbf{Linear Attention Methods.}
    \label{related:linear}
Traditional self-attention~\cite{vaswani2017attention} exhibits quadratic growth in computational complexity as sequence length increases, leading to enormous resource overhead. 
To address this issue, Linear Attention approximates the original dot-product similarity calculation by introducing techniques such as a kernel function \cite{katharopoulos_transformers_2020, choromanski2020rethinking} or by employing low-rank decomposition~\cite{wang2020linformer}. By changing the order of computation or using matrix decomposition, these methods avoid the explicit construction of the attention matrix, which has a quadratic complexity.
This approach significantly reduces the computational and memory burden of Transformers for long-sequence tasks. 
Vision LSTM~\cite{beck:24xlstm}, a linear RNN methods, uses a stacked LSTM structure to extract multi-scale visual features and integrates temporal information through gating mechanisms, achieving strong performance in tasks such as video understanding. Mamba~\cite{gu2023mamba} and its variant Video Mamba~\cite{li2024videomambastatespacemodel}, based on state-space models (SSMs), efficiently model videos and capture rich spatiotemporal dependencies in long-sequence scenarios, demonstrating excellent scalability and expressive power.
However, these methods have not yet been thoroughly validated in the field of medical ultrasound imaging. This domain is particularly challenging due to its complex noise and low contrast. 
Optimizing model structures and training strategies for medical scenarios remains key to further unleashing the potential of linear attention.